% =============================================================================
% m21_ldt -- Statistical analysis (Methods) and Results
% Draft for review. All statistics are taken from the R output of
%   m21_ldt_single_trial.R        (trial-level data set, behavioural scoring)
%   m21_ldt_single_trial_roi.R    (ROI models and ERP reliability)
%   m21_ldt_behaviour.R           (accuracy, RT and behavioural reliability)
%   m21_ldt_single_trial_scalp.R  (supplementary grid analyses)
% Lines marked "% CHECK" need confirming before submission.
% Requires: booktabs, amsmath (for \text in math mode)
% =============================================================================

\subsection{Statistical analysis}

\paragraph{Task performance.}
Behavioural responses were scored from the raw experiment logs. Sensitivity
($d'$) was computed for each participant across all trials, with a log-linear
correction for extreme rates. The three participants excluded for $d' < 0.5$
(see Participants) had $d'$ of 0.10--0.35 and overall accuracy of 52--55\%. In
the analysed sample of 56, $d'$ averaged 2.10 (range 0.73--3.08).

\paragraph{Individual-difference measures.}
Orthographic Sensitivity (OS) and Language Proficiency (LP) were derived from a
principal components analysis of the vocabulary, spelling and author recognition
test scores of the 56 analysed participants. Scores were standardised before the
analysis. The first component (52.7\% of the variance) loaded positively on all
three measures (vocabulary .65, spelling .45, author recognition .62) and was
taken as LP. The second component (28.7\% of the variance) loaded positively on
spelling (.89) and negatively on vocabulary ($-.23$) and author recognition
($-.40$), and was taken as OS. The two scores are uncorrelated by construction
and are centred on zero ($SD_{\text{LP}} = 1.26$, $SD_{\text{OS}} = 0.93$).

\paragraph{Behavioural scoring.}
A response was correct if the key pressed matched the required response.
Responses were collected for 1,000\,ms from target onset; trials with no
response in that window were scored as errors.
Trials on which another key was pressed, or with a recorded RT of 0\,ms, were
treated as missing (about 3 per participant). A capital Z was scored as z.

\paragraph{ERP measures.}
Single-trial ERPs were obtained by assigning each critical item to its own bin.
Mean amplitude was measured in two windows, 200--300\,ms (N250) and
300--500\,ms (N400), relative to a $-200$ to 0\,ms baseline. Analyses used a
central--parietal region of interest (ROI): C3, Cz, C4, CP1, CP2, CP5, CP6, P3,
Pz and P4. P7 and P8 were excluded as too lateral. Amplitude was averaged over
the ROI electrodes on each trial.

\paragraph{Trial and item inclusion.}
ERP and RT analyses used correct trials only; trials rejected for artefacts were
also excluded from the ERP analyses. Four nonwords were removed because they
were presented incorrectly: \textit{trueness} and \textit{boil} are real English
words, \textit{ime} was presented without its base, and \textit{adoptish} was
presented in its complex form in both lists, so it had no simple counterpart.
One word, \textit{design}, was removed because it is not suffixed.

\paragraph{Mixed-effects models.}
All analyses used mixed-effects models in R 4.3.3 with \texttt{lme4} 1.1.35.1
\citep{batesFittingLinearMixedeffects2015}
and \texttt{lmerTest} 3.1.3 \citep{kuznetsovaLmerTestPackageTests2017}, with
Satterthwaite degrees of freedom for linear models. Item properties entered as
continuous predictors: log base frequency, log family size, and, for words, log
whole-word frequency and word length, each standardised across items. Two-level
factors were coded $-0.5$/$+0.5$ (complexity: simple/complex; window: N250/N400).
OS and LP entered as continuous predictors.

\textit{Behaviour.} Accuracy was analysed with logistic models and RT on correct
trials with linear models of log RT. Nonword models included complexity, log
base frequency, log family size, OS and LP; word models included log family
size, log base frequency, log whole-word frequency, length, OS and LP. OS and LP
entered as main effects only, as indices of overall efficiency.

\textit{ERPs.} Three models tested the hypotheses. For words, separate models
were fitted to each window. The N250 model tested whether OS modulates the base
frequency effect (Base Frequency $\times$ OS), with LP, family size, whole-word
frequency and length as covariates. The N400 model tested whether LP modulates
the family size effect (Family Size $\times$ LP), with OS, base frequency,
whole-word frequency and length as covariates. For nonwords, both windows entered
a single model, Complexity $\times$ Window $\times$ (OS $+$ LP), with log base
frequency and log family size (each by window) as covariates. The Complexity
$\times$ OS $\times$ Window interaction tests whether OS modulates the complexity
effect more at the N250 than at the N400.

All models included crossed random effects for participants and items. Each model
began with by-participant and by-item random slopes for the predictors of
interest; terms were removed when the model was singular or did not converge,
until a model converged cleanly \citep{batesParsimoniousMixedModels2018a}.
The final structures are given with each result. When a random slope was removed
because its variance was estimated at zero, the corresponding fixed effect is
tested against trial-level variability, which produces large degrees of freedom.
Interactions were followed up with \texttt{emmeans} 1.10.0
\citep{lenthEmmeansEstimatedMarginal2024}
by estimating the item effect at the mean of the moderator and at $\pm 1\,SD$.
Effect sizes ($\eta^2_p$) were computed from $t$ and its degrees of freedom and
are reported for significant effects in linear models.

Because the words were selected in small and large family size groups, and high
and low base frequency groups, the word ERP models were also run with these
groups in place of the continuous measures. Continuous measures were treated as
primary because dichotomising a continuous variable discards information and
reduces power \citep{maccallumPracticeDichotomizationQuantitative2002}.

\paragraph{Reliability.}
To assess whether participants' effects could be measured reliably enough to
support individual-difference analyses, split-half reliability was estimated
for each participant's experimental effects (complex $-$ simple for nonwords;
large $-$ small family size and high $-$ low base frequency for words) in RT,
accuracy and ERP amplitude, and for each participant's mean ERP amplitude, in
each window \citep{parsonsPsychologicalScienceNeeds2019}.
Each participant's trials were split at random into two halves within each
condition, the effect was computed in each half, the two halves were correlated
across participants and the correlation was corrected with the Spearman--Brown
formula. This was repeated 2,000 times; the median and the 2.5th and 97.5th
percentiles across splits are reported.

\paragraph{Supplementary analyses.}
To check that the ROI did not miss effects with a different scalp distribution,
the ERP analyses were repeated over a $3 \times 3$ grid of electrodes (F3, Fz,
F4, C3, Cz, C4, P3, Pz, P4) with Laterality and Anteriority as factors,
following \citet{yumERPStudyInitial2014}.
These analyses are reported in the Supplementary Material.

% =============================================================================
\section{Results}

\subsection{Behaviour}

\paragraph{Nonwords.}
Complex nonwords were harder to reject than simple nonwords. Accuracy was lower
for complex nonwords (model estimate 65.3\%, 95\% CI [60.0, 70.2]) than for
simple nonwords (88.6\%, [85.4, 91.2]), odds ratio = 0.24, 95\% CI [0.20, 0.29],
$z = -16.62$, $p < .001$, and correct responses were 5.3\% slower, 95\% CI
[3.6, 7.0] (723 vs.\ 687\,ms), $t(79.8) = 6.38$, $p < .001$, $\eta^2_p = .34$
(Table~\ref{tab:behaviour}). Nonwords with higher-frequency bases were also
harder: accuracy was lower, odds ratio = 0.77 per $SD$, $z = -2.86$, $p = .004$,
and RTs were longer, 2.3\% per $SD$, $t(99.6) = 3.81$, $p < .001$,
$\eta^2_p = .13$. Family size had no reliable effect on either measure
($ps \geq .07$).

Participants with higher LP rejected nonwords more accurately, odds ratio = 1.31
per unit, 95\% CI [1.12, 1.53], $z = 3.44$, $p < .001$. OS did not reliably
predict nonword accuracy, odds ratio = 1.21, 95\% CI [0.97, 1.50], $z = 1.71$,
$p = .086$. Neither measure predicted nonword RT (OS: $p = .143$; LP:
$p = .201$).

\paragraph{Words.}
Words with larger families, higher-frequency bases and higher whole-word
frequency were recognised faster (family size: $-1.5\%$ per $SD$, $t(89.9) =
-3.22$, $p = .002$, $\eta^2_p = .10$; base frequency: $-1.3\%$, $t(96.9) = -2.50$,
$p = .014$, $\eta^2_p = .06$; whole-word frequency: $-2.0\%$, $t(91.1) = -4.02$,
$p < .001$, $\eta^2_p = .15$), and longer words more slowly ($+1.5\%$,
$t(91.8) = 3.22$, $p = .002$, $\eta^2_p = .10$). Accuracy was higher for words
with higher whole-word frequency, odds ratio = 1.34, $z = 3.04$, $p = .002$, and
higher base frequency, odds ratio = 1.26, $z = 2.38$, $p = .017$. OS and LP did
not predict word RT or accuracy ($ps \geq .14$).

\begin{table}[htbp]
\centering
\caption{Behavioural models}
\label{tab:behaviour}
\begin{tabular}{lrrlrr}
\toprule
\multicolumn{3}{c}{Accuracy (logistic)} & \multicolumn{3}{c}{Log RT, correct trials} \\
\cmidrule(lr){1-3}\cmidrule(lr){4-6}
Term & OR & $p$ & Term & \% change & $p$ \\
\midrule
\multicolumn{6}{l}{\textit{Nonwords}} \\
Complexity      & 0.24 & $<.001$ & Complexity      & 5.26    & $<.001$ \\
Base frequency  & 0.77 & .004    & Base frequency  & 2.29    & $<.001$ \\
Family size     & 1.02 & .798    & Family size     & $-$1.06 & .071 \\
OS              & 1.21 & .086    & OS              & $-$2.28 & .143 \\
LP              & 1.31 & $<.001$ & LP              & $-$1.47 & .201 \\
\multicolumn{6}{l}{\textit{Words}} \\
Family size     & 1.13 & .161    & Family size     & $-$1.46 & .002 \\
Base frequency  & 1.26 & .017    & Base frequency  & $-$1.29 & .014 \\
Word frequency  & 1.34 & .002    & Word frequency  & $-$1.99 & $<.001$ \\
Length          & 0.95 & .538    & Length          & 1.50    & .002 \\
OS              & 1.10 & .371    & OS              & $-$2.13 & .154 \\
LP              & 1.12 & .146    & LP              & $-$0.58 & .595 \\
\bottomrule
\end{tabular}

\vspace{2pt}
\begin{minipage}{0.85\linewidth}\footnotesize
\textit{Note.} OR = odds ratio. \% change = change in RT, from the exponentiated
log-RT coefficient. Item predictors per $SD$ across items; OS and LP per unit
($SD_{\text{OS}} = 0.93$, $SD_{\text{LP}} = 1.26$); complexity = complex vs.\
simple. Random effects: nonword accuracy, correlated by-participant intercepts
and complexity slopes and by-item intercepts; nonword RT, correlated
by-participant and by-item intercepts and complexity slopes; word models,
by-participant and by-word intercepts.
\end{minipage}
\end{table}

\subsection{Data retained for the ERP analyses}
Participants contributed on average 82.4 correct, artefact-free word trials per
window (range 38--94) and 66.8 nonword trials (range 29--88), of which 28.7
were complex (range 9--45) and 38.1 simple (range 20--48).

\subsection{Reliability}
Participants' mean ERP amplitudes were measured reliably in both windows
(Table~\ref{tab:reliability}). Their ERP effects were not: split-half
reliability was close to or below zero for the complexity, family size and base
frequency effects in both windows. Negative values arise when the variation
between participants' effects is smaller than the variation expected from
measurement error alone, and indicate that the effects carry no reliable
individual differences. The behavioural effects were more reliable than the ERP
effects but still low: .33 for the RT complexity effect, and .12 to .20 for the
others.

\begin{table}[htbp]
\centering
\caption{Split-half reliability (Spearman--Brown corrected)}
\label{tab:reliability}
\begin{tabular}{llccc}
\toprule
Stimuli & Measure & N400 & N250 & Behaviour \\
\midrule
Words    & Mean amplitude        & .87 [.80, .92]          & .85 [.77, .90]          & \\
         & Family size effect    & $-$.04 [$-$.64, .33]    & .08 [$-$.42, .40]       & RT: .15 [$-$.22, .44] \\
         &                       &                         &                         & Acc.: .20 [$-$.16, .47] \\
         & Base frequency effect & $-$.38 [$-$1.14, .12]   & $-$.65 [$-$1.51, $-$.04] & \\
Nonwords & Mean amplitude        & .86 [.79, .91]          & .78 [.68, .86]          & \\
         & Complexity effect     & $-$.70 [$-$1.57, $-$.06] & $-$.11 [$-$.66, .30]    & RT: .33 [.00, .58] \\
         &                       &                         &                         & Acc.: .12 [$-$.28, .43] \\
\bottomrule
\end{tabular}

\vspace{2pt}
\begin{minipage}{0.9\linewidth}\footnotesize
\textit{Note.} Values are medians across 2,000 random splits; brackets give the
2.5th and 97.5th percentiles across splits. ERP values are for the ROI. RT
reliability uses correct trials.
\end{minipage}
\end{table}

\subsection{ERPs: words}

\paragraph{N400: family size and Language Proficiency.}
Words with larger families elicited more positive N400 amplitudes, $b = 0.32$\,$\mu$V
per $SD$ of log family size, 95\% CI [0.03, 0.61], $t(91.9) = 2.20$, $p = .030$,
$\eta^2_p = .05$. LP did not modulate this effect, $b = -0.02$, 95\% CI
[$-0.22$, 0.18], $t(4480.1) = -0.24$, $p = .811$ (Table~\ref{tab:words}). The
family size effect was similar at low LP ($-1\,SD$: 0.35\,$\mu$V, $p = .072$)
and high LP ($+1\,SD$: 0.29\,$\mu$V, $p = .139$). Per $SD$ of LP, the 95\% CI
for the interaction was $-0.28$ to $0.22$\,$\mu$V.

Higher-frequency words also elicited more positive N400 amplitudes, $b = 0.45$,
95\% CI [0.13, 0.76], $t(93.5) = 2.83$, $p = .006$, $\eta^2_p = .08$, and
participants with higher OS showed more positive N400 amplitudes overall,
$b = 1.13$ per unit OS, 95\% CI [0.38, 1.88], $t(53.4) = 3.01$, $p = .004$,
$\eta^2_p = .14$. Base frequency, length and LP had no effect ($ps \geq .78$).
The final random-effects structure comprised by-participant and by-word
intercepts.

With family size coded as small versus large, the family size effect was in the
same direction but not significant, $b = 0.53$\,$\mu$V, 95\% CI [$-0.04$, 1.11],
$t(93.2) = 1.84$, $p = .068$; the interaction with LP remained null, $p = .827$.

\paragraph{N250: base frequency and Orthographic Sensitivity.}
Base frequency had no effect at the N250, $b = 0.03$, 95\% CI [$-0.27$, 0.32],
$t(4555.3) = 0.18$, $p = .857$, and OS did not modulate it, $b = -0.07$, 95\% CI
[$-0.35$, 0.22], $t(4554.6) = -0.46$, $p = .647$ (Table~\ref{tab:words}). The
base frequency effect was near zero at low OS ($0.09$\,$\mu$V, $p = .652$) and
high OS ($-0.03$\,$\mu$V, $p = .868$). Per $SD$ of OS, the 95\% CI for the
interaction was $-0.33$ to $0.20$\,$\mu$V. The high versus low coding gave the
same result, interaction $p = .640$. The variance between words, and between
participants' base frequency effects, was estimated at zero, so the final model
included by-participant intercepts only.

Family size, a covariate in this model, was associated with more positive N250
amplitudes, $b = 0.32$, 95\% CI [0.07, 0.58], $t(4553.6) = 2.52$, $p = .012$,
$\eta^2_p < .01$. No other predictor was significant ($ps \geq .09$).

\begin{table}[htbp]
\centering
\caption{Single-trial ERP models for words (ROI)}
\label{tab:words}
\begin{tabular}{lrrrr@{\hspace{2em}}lrrrr}
\toprule
\multicolumn{5}{c}{N400} & \multicolumn{5}{c}{N250} \\
\cmidrule(lr){1-5}\cmidrule(lr){6-10}
Term & $b$ & $SE$ & $t$ & $p$ & Term & $b$ & $SE$ & $t$ & $p$ \\
\midrule
Family size (FS)  &  0.32 & 0.14 &  2.20 & .030 & Base frequency (BF) &  0.03 & 0.15 &  0.18 & .857 \\
LP                & $-$0.01 & 0.28 & $-$0.04 & .969 & OS                  &  0.34 & 0.37 &  0.94 & .353 \\
FS $\times$ LP    & $-$0.02 & 0.10 & $-$0.24 & .811 & BF $\times$ OS      & $-$0.07 & 0.14 & $-$0.46 & .647 \\
OS                &  1.13 & 0.37 &  3.01 & .004 & LP                  & $-$0.12 & 0.27 & $-$0.43 & .668 \\
Base frequency    &  0.05 & 0.17 &  0.27 & .784 & Family size         &  0.32 & 0.13 &  2.52 & .012 \\
Word frequency    &  0.45 & 0.16 &  2.83 & .006 & Word frequency      &  0.09 & 0.14 &  0.65 & .513 \\
Length            &  0.01 & 0.15 &  0.05 & .958 & Length              & $-$0.22 & 0.13 & $-$1.67 & .094 \\
\bottomrule
\end{tabular}

\vspace{2pt}
\begin{minipage}{0.95\linewidth}\footnotesize
\textit{Note.} Estimates in $\mu$V. Item predictors are standardised across words;
OS and LP are component scores ($SD_{\text{OS}} = 0.93$, $SD_{\text{LP}} = 1.26$).
Degrees of freedom (Satterthwaite) are given in the text.
\end{minipage}
\end{table}

\subsection{ERPs: nonwords}
Complex and simple nonwords did not differ at the N250, $-0.10$\,$\mu$V, 95\% CI
[$-0.83$, 0.63], $t(143) = -0.27$, $p = .788$, or at the N400, 0.14\,$\mu$V, 95\%
CI [$-0.59$, 0.87], $t(143) = 0.37$, $p = .709$. OS did not modulate the
complexity effect differently in the two windows, Complexity $\times$ OS $\times$
Window, $b = 0.26$, 95\% CI [$-0.63$, 1.14], $t(7144.9) = 0.56$, $p = .573$
(Table~\ref{tab:nonwords}), and did not modulate it in either window (N250:
$b = -0.21$, 95\% CI [$-0.86$, 0.44], $t(173) = -0.63$, $p = .528$; N400:
$b = 0.05$, 95\% CI [$-0.61$, 0.70], $t(173) = 0.14$, $p = .888$). LP did not
modulate the complexity effect either ($ps \geq .76$). The final random-effects
structure comprised uncorrelated by-participant intercepts and complexity and
window slopes, and uncorrelated by-item intercepts and complexity slopes.

No other term involving OS or LP was significant (Table~\ref{tab:nonwords}).

\begin{table}[htbp]
\centering
\caption{Single-trial ERP model for nonwords (ROI)}
\label{tab:nonwords}
\begin{tabular}{lrrrrr}
\toprule
Term & $b$ & $SE$ & $df$ & $t$ & $p$ \\
\midrule
Complexity (C)                    &  0.02 & 0.31 &   67.6 &  0.06 & .950 \\
Window (W)                        &  0.56 & 0.35 &   56.7 &  1.63 & .109 \\
OS                                &  0.63 & 0.35 &   54.6 &  1.81 & .075 \\
LP                                &  0.00 & 0.26 &   55.0 &  0.00 & .999 \\
C $\times$ W                      &  0.22 & 0.42 & 7145.8 &  0.53 & .599 \\
C $\times$ OS                     & $-$0.08 & 0.24 &   48.6 & $-$0.34 & .737 \\
C $\times$ LP                     &  0.03 & 0.18 &   50.7 &  0.17 & .868 \\
W $\times$ OS                     &  0.63 & 0.38 &   57.0 &  1.67 & .100 \\
W $\times$ LP                     &  0.15 & 0.28 &   59.1 &  0.54 & .591 \\
C $\times$ W $\times$ OS          &  0.26 & 0.45 & 7144.9 &  0.56 & .573 \\
C $\times$ W $\times$ LP          &  0.09 & 0.34 & 7154.5 &  0.26 & .792 \\
\bottomrule
\end{tabular}

\vspace{2pt}
\begin{minipage}{0.75\linewidth}\footnotesize
\textit{Note.} Estimates in $\mu$V. Complexity and Window coded $-0.5$/$+0.5$
(simple/complex; N250/N400). Log base frequency and log family size (and their
interactions with Window) were included as covariates and were not significant
($ps \geq .64$).
\end{minipage}
\end{table}

\subsection{Supplementary grid analyses}
Over the $3 \times 3$ grid, no term combining complexity with OS was significant,
alone or with Window, Laterality or Anteriority ($ps \geq .13$). Complexity
interacted with Laterality, $F(2, 59053) = 3.76$, $p = .023$, but the complexity
effect was not significant at any laterality ($ps \geq .55$). Family size
interacted with Anteriority, $F(2, 72907) = 7.53$, $p < .001$: the family size
effect increased from frontal (0.26\,$\mu$V, $p = .383$) to central
(0.52\,$\mu$V, $p = .082$) and parietal sites (0.61\,$\mu$V, $p = .044$), and
was significant only at parietal sites. These scalp interactions are tested against
electrode-level variability within trials and should be interpreted cautiously.
Full results are in the Supplementary Material.
% CHECK: describe the supplementary models fully in the Supplement (nine grid
% electrodes on each trial; small/large and high/low groups for words).
